% Generated draft: requires mathematical and source review.
\documentclass[11pt]{article}
\usepackage[margin=25mm]{geometry}
\usepackage{fontspec}
\setmainfont{Latin Modern Roman}
\usepackage{amsmath,amssymb,mathtools,unicode-math}
\setmathfont{Latin Modern Math}
\usepackage{enumitem,needspace,xurl,hyperref}
\setlength{\parindent}{0pt}
\setlength{\parskip}{5pt}
\setlength{\emergencystretch}{3em}
\newcommand{\N}{\mathbb{N}}
\newcommand{\Z}{\mathbb{Z}}
\newcommand{\Q}{\mathbb{Q}}
\newcommand{\R}{\mathbb{R}}
\newcommand{\C}{\mathbb{C}}
\newcommand{\F}{\mathbb{F}}
\newcommand{\A}{\mathbb{A}}
\newcommand{\PP}{\mathbb{P}}
\newcommand{\E}{\mathbb{E}}
\newcommand{\eps}{\varepsilon}
\newcommand{\dd}{\,\mathrm{d}}
\begin{document}
\section*{UnsolvedMath ID: 30006283}
\textbf{Source problem code:} catalogue entry
\subsection{Definitions and mathematical statement}
Write \(X=SU_C(2,\Lambda)\) for the rank-two bundle moduli space with fixed determinant \(\Lambda\) named in the supplied record. The record does not specify the hypotheses on \(C\), the degree of \(\Lambda\), the ground field, or the stability convention defining this space. These must be recovered from the cited report before a fully specified conjecture can be stated.

A semiorthogonal decomposition of a triangulated category \(\mathcal D\) is an ordered collection of full triangulated subcategories \(\mathcal A_1,\ldots,\mathcal A_m\) that generate \(\mathcal D\) and satisfy
\[
\operatorname{Hom}_{\mathcal D}(A_j,A_i[k])=0
\quad\text{for every }j>i,\ k\in\Z,
\ A_j\in\mathcal A_j,\ A_i\in\mathcal A_i.
\]
The source refers specifically to the semiorthogonal decomposition of \(X\) in its Theorem~1, but neither that theorem nor its ambient category and components are supplied. Interpreting the ambient category as \(D^b(\operatorname{Coh}(X))\) would require confirmation from that theorem.

For the geometric hypotheses and quantum-cohomology parameters intended by the source, let
\[
H=H^*(X,\C),\qquad T(\alpha)=c_1\star\alpha,
\]
where \(\star\) denotes the intended quantum product and \(c_1\) the class appearing in the source. Its usual interpretation as \(c_1(T_X)\), as well as the quantum product and any specialization needed to obtain a complex-linear endomorphism of \(H\), requires confirmation. If these data give an endomorphism of finite-dimensional \(H\), its generalized eigenspaces are
\[
H_\mu=\ker\bigl((T-\mu\operatorname{id}_H)^N\bigr)
\quad(N\text{ sufficiently large}),\qquad
H=\bigoplus_{\mu\in\operatorname{Spec}(T)}H_\mu.
\]
The original statement calls this decomposition orthogonal, without specifying the pairing.

\textbf{Source Conjecture 5, with unresolved specifications.} The particular semiorthogonal decomposition of \(SU_C(2,\Lambda)\) given in Theorem~1 of the cited report is atomic: it is compatible with the decomposition of \(H^*(SU_C(2,\Lambda),\C)\) into generalized eigenspaces of quantum multiplication by \(c_1\), in precisely the sense of the Sanda--Shamoto conjecture cited there as reference~[13]. The supplied record does not reproduce that compatibility condition. Thus this is a reference-dependent statement of the intended conjecture, and the missing theorem and compatibility definition prevent a complete intrinsic formulation.

The original extraction also contains the fragment that the eigenspace corresponding to \(8\lambda\) is isomorphic to
\[
H^*\bigl(\operatorname{Sym}^{g-1-|\lambda|}(C),\C\bigr).
\]
For a nonnegative integer \(n\), \(\operatorname{Sym}^n(C)=C^n/\mathfrak S_n\). The fragment does not specify the range or meaning of \(\lambda\), define \(g\), identify which operator has eigenvalue \(8\lambda\), or specify the structure preserved by the isomorphism. Its placement after the conjecture appears disrupted. It is retained as unresolved source context, not asserted as an independently specified theorem or an additional conjectural conclusion.
\subsection{Short English statement}
The source conjectures that its specific semiorthogonal decomposition of the rank-two fixed-determinant bundle moduli space matches the generalized quantum-cohomology eigenspace decomposition in the Sanda--Shamoto sense. The geometric hypotheses, proposed categorical components, quantum parameters, and precise matching condition are missing from the supplied extract, so the statement remains reference-dependent and requires mathematical review.
\subsection{Sources}
\begin{itemize}
\item Local UnsolvedMath registry, permanent ID \(30006283\), title ``Atomic Decomposition of Rank-Two Bundle Moduli,'' source code \texttt{catalogue entry}, in the Oberwolfach Reports Open Problems collection. The supplied extract identifies the target as Conjecture~5 and refers to Theorem~1. No individual proposer is supplied; the recorded proposed year is 2025.
\item \emph{Recent Trends in Algebraic Geometry} (2025), Oberwolfach Reports, pp.~1451--1516, DOI \texttt{10.4171/owr/2025/28}. Supplied source report and literature-source link: \url{https://doi.org/10.4171/owr/2025/28}. The supplied background describes Oberwolfach Reports as documenting workshops and their open-problem discussions.
\item Sanda--Shamoto, ``analogue of Dubrovin's conjecture,'' as labeled in the supplied literature sources: \url{https://arxiv.org/abs/1705.05989}. The original extract attributes its precise compatibility condition to the Sanda--Shamoto conjecture, reference~[13] in the report; that definition is not reproduced here.
\item Supplied literature assessment dated exactly \texttt{2026-08-22}: the record labels the problem open and reports that the claimed atomic decomposition remains conjectural, with no general proof found by its cutoff date. This is supplied dated curation derived from the OWR material, not an independently verified literature assessment. The background requests independent verification of time-sensitive status and cited claims.
\item The registry marks the statement \texttt{corrected\_verified} and reports that the official report was checked, the question isolated from surrounding exposition, and notation repaired. These are supplied provenance claims; no source visit or independent verification is asserted in this draft.
\end{itemize}
\end{document}
